\section{Conclusão}
\label{sec:conclusao}

Os três blocos foram projetados a partir das funções derivadas da
matrícula 21101175, com as deduções apresentadas na
Seção~\ref{sec:projeto}: cascata de somadores inversores nos blocos A
e B e amplificador de instrumentação com ganho $12{,}5$ no bloco C. O
circuito completo utiliza sete LM741 alimentados em
$\pm$\SI{15}{\volt}.

A simulação confirmou o projeto em três níveis. Os ganhos individuais
de cada entrada reproduziram os coeficientes teóricos com desvio
inferior a $0{,}1\%$. Cada bloco isolado reproduziu a respectiva
função com erro máximo abaixo de \SI{0.3}{\milli\volt}. A resposta
global acompanhou a referência ideal com erro máximo de
$9{,}3\times10^{-4}$~V em uma excursão de \SI{11.4}{\volt}. Com o
macromodelo comercial do LM741 o erro máximo subiu para $0{,}27$~V
(cerca de 2{,}4\% da excursão), efeito da queda do ganho de malha
aberta com a frequência e do offset de entrada do componente real, sem
alterar o comportamento projetado. O ensaio adicional com um
amplificador moderno (ADA4610), na mesma alimentação, reduziu o erro
para menos de 0{,}1\% da excursão, confirmando que o desvio se deve às
limitações do LM741 e não à topologia.

Como continuidade do trabalho, sugere-se elevar a frequência dos
estímulos até evidenciar as limitações de banda e de taxa de variação
(\emph{slew rate}) do LM741 real, e avaliar o efeito da tolerância dos
resistores em uma montagem física.
